\chapter*{Introducción}
\addcontentsline{toc}{chapter}{Introducción}

% La introducción no cuenta como primer capítulo

\noindent Desde 1998 la Asociación Mexicana de Agencias de Inteligencia de Mercado y Opinión A.C. (AMAI) realiza el \emph{Estudio Anual de la Industria de Investigación de Mercados y Opinión Pública en México} (Estudio Anual de la AMAI), dirigido a sus asociados, las cuales son organizaciones de inteligencia de mercado y opinión que realizan sus actividades en México.

En mayo de 2020, la AMAI encargó al Centro de Ciencia de Datos ITAM el procesamiento de la información generada en la \emph{Edición XXII (2019-2020)} de dicho estudio, con el propósito de consolidar documentos de análisis que faciliten su entendimiento y ayuden a encontrar elementos de valor para la industria.

Bajo dicho contexto, el presente trabajo expone los principios metodológicos considerados, como asistente de investigación del Centro de Ciencia de Datos del ITAM, para efectuar el correspondiente procesamiento y análisis de la información recibida en la \emph{Edición XXII (2019-2020)} del Estudio Anual de la AMAI, a través del desarrolló de un producto de datos encaminado hacia dicho propósito.

El trabajo se ha estructurado a manera de capítulos, donde el primero de ellos aborda el contexto de la asociación AMAI y de los objetivos que se buscan alcanzar a través  de la \emph{Edición XXII (2019-2020)} de este estudio. Por otra parte, el segundo capítulo plantea dichos objetivos como problemas  a  resolverse desde  la  óptica  de  ciencia  de  datos. En este sentido, el tercer capítulo describe el producto de datos diseñado para abordar los problemas previos, planteando un proceso de automatización para consolidar reportes con la información de coyuntura que es de interés para la industria de inteligencia de negocios e investigación de mercado, desde el punto de vista técnico y conceptual. Finalmente, se ofrece en un cuarto capítulo que aborda la versión pública de los resultados obtenidos con dicho proceso para la \emph{Edición XXII (2019-2020)} del Estudio Anual de la AMAI, cerrando con un apartado donde se exponen las principales conclusiones y recomendaciones.

% Se sugiere que el primer párrafo de cada sección no tenga sangría: \noindent
